% Source: https://artofproblemsolving.com/community/c3h2794859_2017_geometry_bee_invitational_round
% Stable identity: aops-topic-2794859-q10
The angle-bisector theorem gives $NB:NC=5:7$. Triangles $PNB,PNC$ have supplementary angles at $N$, so the sine rule yields $\dfrac{NB}{NC}=\dfrac{PB\sin\angle BPN}{PC\sin\angle CPN}$. Since $P,N,M$ are collinear and chords $BM,CM$ are equal, the two sine factors are equal (their angles are equal or supplementary). Hence $PB:PC=5:7$, and $\boxed{PC=\dfrac{21}{5}}$.
\begin{asy}
usepackage("amsmath");
size(7cm);
pair B=(0,0),C=(8,0),A=(2.5,5*sqrt(3)/2),O=(4,1/sqrt(3)),M=(4,8/sqrt(3)),Nn=(10/3,0),v=Nn-M;real t=-2*dot(v,M-O)/dot(v,v);pair P=M+t*v;draw(circle(O,7/sqrt(3)));draw(A--B--C--cycle);draw(M--P);draw(A--Nn,dashed);label("A",A,NW);label("B",B,SW);label("C",C,SE);label("M",M,N);label("N",Nn,S);label("P",P,SE);label("schematic",(4,-4));
\end{asy}
